% ═══════════════════════════════════════════════════════════════════════
\section{Au-delà de trois dimensions}
% ═══════════════════════════════════════════════════════════════════════

\begin{frame}{Le décathlon}
  \small
  \begin{itemize}
    \item Dix épreuves sur deux jours : quatre courses ($100$ m, $400$ m,
      $110$ m haies, $1500$ m), trois sauts (longueur, hauteur, perche),
      trois lancers (poids, disque, javelot).
    \item $41$ résultats de 2004 : Décastar de Talence et Jeux olympiques
      d'Athènes. Chaque résultat est un point de $\R^{10}$ ; aucun dessin
      direct n'est possible.
  \end{itemize}
  \scriptsize Tableau 5.3 --- Quelques décathloniens en 2004 (secondes pour les
  courses, mètres pour les sauts et les lancers).
  \begin{center}
    \scriptsize
    \setlength{\tabcolsep}{4pt}
    \begin{tabular}{lcccccccccc}
      \toprule
      Athlète & $100$ & Long. & Poids & Haut. & $400$ & Haies & Disque & Perche & Jav. & $1500$ \\
      \midrule
      Sebrle & $\num{10.85}$ & $\num{7.84}$ & $\num{16.36}$ & $\num{2.12}$ & $\num{48.36}$ & $\num{14.05}$ & $\num{48.72}$ & $\num{5.00}$ & $\num{70.52}$ & $\num{280.01}$ \\
      Karpov & $\num{10.50}$ & $\num{7.81}$ & $\num{15.93}$ & $\num{2.09}$ & $\num{46.81}$ & $\num{13.97}$ & $\num{51.65}$ & $\num{4.60}$ & $\num{55.54}$ & $\num{278.11}$ \\
      Drews & $\num{10.87}$ & $\num{7.38}$ & $\num{13.07}$ & $\num{1.88}$ & $\num{48.51}$ & $\num{14.01}$ & $\num{40.11}$ & $\num{5.00}$ & $\num{51.53}$ & $\num{274.21}$ \\
      Casarsa & $\num{11.36}$ & $\num{6.68}$ & $\num{14.92}$ & $\num{1.94}$ & $\num{53.20}$ & $\num{15.39}$ & $\num{48.66}$ & $\num{4.40}$ & $\num{58.62}$ & $\num{296.12}$ \\
      \midrule
      moyenne & $\num{11.00}$ & $\num{7.26}$ & $\num{14.48}$ & $\num{1.98}$ & $\num{49.62}$ & $\num{14.61}$ & $\num{44.33}$ & $\num{4.76}$ & $\num{58.32}$ & $\num{279.02}$ \\
      \bottomrule
    \end{tabular}
  \end{center}
  \small Unités et dispersions différentes : on réduit ; la variance totale
  vaut $10$.
\end{frame}

\begin{frame}[fragile]{Programme 5.3 --- L'ACP du décathlon}
\begin{lstlisting}
deca = pd.read_csv("decathlon.csv", index_col=0)
X = deca.iloc[:, :10].to_numpy()        # les dix épreuves
valeurs, vecteurs, Y = acp(X)
print("valeurs propres :", valeurs.round(2))
print("% cumulés :", (100 * np.cumsum(valeurs) / valeurs.sum()).round(1))
r = np.corrcoef(Y[:, 0], deca["Points"])[0, 1]
print("corrélation y1 / points :", round(r, 3))
\end{lstlisting}
  \vspace{4pt}
  \small Sortie (la fonction \texttt{acp} est celle du programme 5.1) :
  \begin{center}
    \footnotesize
    \begin{tabular}{lcccccccccc}
      \toprule
      $k$ & $1$ & $2$ & $3$ & $4$ & $5$ & $6$ & $7$ & $8$ & $9$ & $10$ \\
      \midrule
      $\lambda_k$ & $\num{3.27}$ & $\num{1.74}$ & $\num{1.40}$ & $\num{1.06}$ & $\num{0.68}$ & $\num{0.60}$ & $\num{0.45}$ & $\num{0.40}$ & $\num{0.21}$ & $\num{0.18}$ \\
      \% cumulé & $\num{32.7}$ & $\num{50.1}$ & $\num{64.1}$ & $\num{74.7}$ & $\num{81.6}$ & $\num{87.5}$ & $\num{92.1}$ & $\num{96.0}$ & $\num{98.2}$ & $100$ \\
      \bottomrule
    \end{tabular}
  \end{center}
  Corrélation entre $y_1$ et le total officiel des points : $\num{0.956}$.
\end{frame}

\begin{frame}{Figure 5.9 --- Le diagramme des éboulis}
  \begin{columns}[c]
    \begin{column}{0.55\textwidth}
      \FigLivre[0.62\textheight]{fig-acp-eboulis}
      \scriptsize La droite pointillée : le seuil du critère de Kaiser.
    \end{column}
    \begin{column}{0.43\textwidth}
      \small
      \begin{itemize}
        \item Beaucoup moins tranché que pour les manchots : le plan
          principal ne conserve qu'environ la moitié de l'information
          ($\SI{32.7}{\percent}+\SI{17.4}{\percent}$).
        \item Kaiser retiendrait quatre axes ; il en faudrait cinq pour
          dépasser $\SI{80}{\percent}$.
        \item Pas un échec : une information sur les données. Les dix
          épreuves sont assez indépendantes ; on ne résume pas un
          décathlonien par deux nombres.
      \end{itemize}
    \end{column}
  \end{columns}
\end{frame}

\begin{frame}{Lire le plan principal}
  \begin{columns}[c]
    \begin{column}{0.4\textwidth}
      \scriptsize Tableau 5.4 --- Les deux premiers vecteurs propres.
      \begin{center}
        \scriptsize
        \begin{tabular}{lcc}
          \toprule
          Épreuve & $\vec v_1$ & $\vec v_2$ \\
          \midrule
          $100$ m & $-\num{0.428}$ & $\num{0.142}$ \\
          saut en longueur & $\num{0.410}$ & $-\num{0.262}$ \\
          lancer du poids & $\num{0.344}$ & $\num{0.454}$ \\
          saut en hauteur & $\num{0.316}$ & $\num{0.266}$ \\
          $400$ m & $-\num{0.376}$ & $\num{0.432}$ \\
          $110$ m haies & $-\num{0.413}$ & $\num{0.174}$ \\
          lancer du disque & $\num{0.305}$ & $\num{0.460}$ \\
          saut à la perche & $\num{0.028}$ & $-\num{0.137}$ \\
          lancer du javelot & $\num{0.153}$ & $\num{0.241}$ \\
          $1500$ m & $-\num{0.032}$ & $\num{0.360}$ \\
          \bottomrule
        \end{tabular}
      \end{center}
    \end{column}
    \begin{column}{0.58\textwidth}
      \small
      \begin{itemize}
        \item Premier axe : vitesse ($100$ m, haies, $400$ m) négative, saut
          en longueur positif. Un bon temps est un temps court : à droite,
          les athlètes qui courent vite et sautent loin. La performance
          générale.
        \item Second axe : les lancers (poids, disque), mais aussi le $400$ m
          et le $1500$ m. En haut : puissants, mais relativement lents sur
          les courses longues (Casarsa : disque à $\num{48.66}$ m,
          $400$ m en $\num{53.20}$ s).
        \item La perche, aux deux composantes faibles, est mal représentée :
          c'est le troisième axe qui la décrit.
      \end{itemize}
    \end{column}
  \end{columns}
\end{frame}

\begin{frame}{Figure 5.10 --- Les athlètes dans le plan principal}
  \begin{columns}[c]
    \begin{column}{0.58\textwidth}
      \FigLivre[0.72\textheight]{fig-acp-decathlon}
    \end{column}
    \begin{column}{0.4\textwidth}
      \small
      \begin{itemize}
        \item La première composante a été calculée sans le barème officiel,
          qui convertit chaque performance en points selon des tables
          complexes.
        \item Pourtant, sa corrélation avec le total des points vaut
          $\num{0.956}$ : les trois médaillés d'Athènes, Sebrle, Clay et
          Karpov, sont les trois points les plus à droite.
        \item L'ACP a retrouvé seule ce que signifie « être un bon
          décathlonien ».
      \end{itemize}
    \end{column}
  \end{columns}
\end{frame}

\begin{frame}{D'autres applications}
  \begin{itemize}
    \item Génétique : Novembre et ses collègues (2008), près de
      $\num{1400}$ Européens décrits par environ $\num{200000}$ marqueurs.
      Dans le plan principal, convenablement tourné, le nuage reproduit la
      carte de l'Europe.
    \item Images : une photo de $100\times100$ pixels est un point de
      $\R^{\num{10000}}$. Les vecteurs propres d'un grand nombre de visages
      sont des « visages propres » (eigenfaces) ; quelques dizaines de
      coordonnées suffisent à bien décrire un visage.
    \item Compresser des données, éliminer du bruit, réduire le nombre de
      variables avant une autre méthode d'apprentissage automatique.
    \item Souvent la première chose que l'on fait devant un nouveau jeu de
      données.
  \end{itemize}
\end{frame}

\begin{frame}{Les limites de la méthode}
  \begin{itemize}
    \item L'ACP est linéaire : elle cherche des droites et des plans, et
      passe à côté d'une structure courbe (un nuage le long d'un arc de
      cercle).
    \item Elle est sensible aux unités (remarque 5.4) et aux valeurs
      aberrantes, qui peuvent accaparer un axe à elles seules :
      \begin{itemize}
        \item deux élèves atypiques ajoutés aux six font passer la
          corrélation de $\num{0.8}$ à environ $-\num{0.08}$ et basculer
          l'axe principal (exercice 5.5) ;
        \item une masse de manchot saisie à $\num{42000}$ g au lieu de
          $\num{4200}$ g s'empare du deuxième axe, même après réduction
          (exercice 5.10).
      \end{itemize}
      Avant d'interpréter : regarder les données et repérer les individus
      atypiques.
    \item Une grande variance n'est pas toujours une information
      intéressante : la direction qui sépare deux groupes peut être une
      direction de faible variance.
  \end{itemize}
\end{frame}

\begin{frame}{Pour la classe 5.4 --- Des données réelles en classe}
  \begin{Classe}{Des données réelles en classe}
    \small
    \begin{itemize}
      \item Le programme insiste sur l'interprétation de données et le regard
        critique sur l'information statistique.
      \item Sans algèbre linéaire : distribuer les mesures d'une cinquantaine
        de manchots, sans l'espèce. Quel nuage de points, pour quelle paire
        de variables, montre le mieux qu'il y a « plusieurs sortes » de
        manchots ?
      \item L'idée centrale du chapitre : un même ensemble de données,
        regardé sous différents angles, peut révéler ou cacher sa structure.
      \item Élèves plus avancés : un tableur pour calculer les cotes $z$, et
        constater que les variables réduites deviennent comparables.
      \item Construire un « score » combinant plusieurs épreuves d'éducation
        physique, et le comparer à une simple somme des cotes $z$.
    \end{itemize}
  \end{Classe}
\end{frame}
